\section{Q4: Explain the role and importance of model serving and CI/CD in today’s data 
engineering world. Write down the names and briefly explain four frameworks (two 
for model serving and two for CI/CD) that offer model serving and CI/CD 
capabilities.}

\subsection*{Role and Importance}

Training a model is only the first step. To be useful, the model must run as a live service so other applications can send data and receive predictions --- this is model serving. It manages model versions, routes requests, and scales with demand. CI/CD (Continuous Integration and Continuous Deployment) automates the steps between updating a model and deploying it to production. Because data changes over time, models need regular retraining and redeployment. 
Without CI/CD, this is slow and error-prone. Together, model serving and CI/CD let teams update models quickly, catch errors through automated testing, and keep predictions stable in production~\cite{qian2020orchestrating, verbraeken2020survey}.

\subsection*{Four Frameworks}

\begin{enumerate}
  \item \textbf{TensorFlow Serving (Model Serving):} An open-source tool by Google for deploying TensorFlow models. It serves predictions via REST or gRPC APIs and supports multiple model versions at once, which is useful for A/B testing. It is built for high-speed, high-load environments.

  \item \textbf{MLflow Model Serving (Model Serving):} An open-source platform that works with any ML framework. A single command deploys a trained model as a REST API. It also provides a model registry to track versions and safely move models from testing to production.

  \item \textbf{Jenkins (CI/CD):} One of the most widely used CI/CD tools. It can automatically retrain a model when new data arrives, run validation tests, and deploy passing models. Its many plugins connect it to cloud platforms, Docker, and ML frameworks.

  \item \textbf{PerfCI (CI/CD):} A CI/CD framework designed specifically for ML pipelines. It automatically benchmarks model performance such as speed and accuracy. Each time a change is made, it flags regressions before the model reaches production. This makes it especially useful for teams that need to ensure model quality does not drop between updates.
\end{enumerate}
